\documentclass[../../../include/open-logic-section]{subfiles}

\begin{document}

\olfileid{his}{set}{mythology}

\olsection{تاریخ توں ودھ متھ؟}

ہُن اک صدی توں وی ودھ پچھوں ایہناں واقعات ول مڑ کے ویکھدیاں، سانوں ایہ فیصلے سُن کے ای منّ نہیں لینے چاہیدے۔ نتیجے سچ مُچ نویں، دلچسپ تے حیرانی والے سن۔ پر ایہ کِنّے کُ حیران کر دین والے سن؟ تے کی ایہناں سچ مُچ وکھایا کہ سانوں ہندسی سہج سمجھ اُتے بھروسا نہیں کرنا چاہیدا؟

حیرانی والی گل بارے \citet{Gouvea2011} دسدا اے کہ ڈیڈیکائنڈ نوں کانتور دی مشہور گل، «\emph{je le vois, mais je ne le crois
pas}»، اپنے اصل سیاق توں کافی وکھ کر کے پیش کیتی جاندی اے۔ اوہ سیاق ہور کھولھ کے ایویں اے (\citeauthor{Gouvea2011} توں لیا گیا):
\begin{quote}
جے میں تہاڈی مہربانی تے صبر توں اینا کم لینا آں، تاں ایس موضوع لئی میرا شوق معاف کرنا۔ میں جیہڑیاں گلّاں ہُنے تہاڈے ول بھیجیاں، اوہ میرے اپنے لئی وی اینیاں اچانک تے اینیاں نویاں نیں کہ، میرے محترم سجن، جدوں تک تہاڈے کولوں اوہناں دے ٹھیک ہون بارے فیصلہ نہیں ملدا، مینوں چین نہیں آ سکدا۔ جدوں تک تسی میرے نال سہمت نہیں ہوندے، میں بس ایہ آکھ سکدا آں:
\emph{je le vois, mais je ne le crois pas.}
\end{quote}
کانتور جاندا سی کہ اوہدا نتیجہ «اینا اچانک، اینا نواں» سی۔ پر ایہ شک والی گل اے کہ اوہنوں اپنا نتیجہ کدے \emph{ناقابلِ یقین} لگیا ہووے۔ جیویں \citeauthor{Gouvea2011} دسدا اے، کانتور ڈیڈیکائنڈ نوں بس اپنا پیش کیتا ثبوت پرکھن لئی آکھ رِہا سی۔

ہندسی سہج سمجھ دی گل ایہ اے کہ پیانو نے اپنی ساری تھاں بھرن والی منحنی کسے خاکے توں بناں چھپوائی۔ پر ہلبرٹ نے اپنی منحنی چھپواؤن ویلے اپنا مقصد سمجھایا۔ جے قاری «اگّے دِتّی ہندسی سہج سمجھ دی مدد لین»، تاں اوہناں کول پیانو دے نتیجے نوں سمجھن دا صاف طریقہ ہووے گا۔ فیر اوہنے \olref[his][set][pathology]{sec} وچلے ورگے \emph{خاکیاں} دی اک لڑی دِتّی۔

ہور عام گل ایہ اے کہ بھاویں چھپے ہوئے ثبوتاں وچ خاکیاں دا رواج کافی گھٹ گیا اے، ایہ حقیقت نہیں ٹل سکدی کہ ریاضی دان ثبوت دیندیاں \emph{اکثر} خاکے ورتدے نیں۔ (موٹے طور اُتے آکھیے، چنگے ریاضی دان جاندے نیں کہ ہندسی سہج سمجھ اُتے کدوں بھروسا کیتا جا سکدا اے۔)

مُکدی گل ایہ اے کہ رولا سُن کے ای یقین نہ کرو، یا گھٹ توں گھٹ ایہنوں بناں پرکھے نہ منّو۔ ہور پڑھنا ہووے تاں \citet{Giaquinto2007} ویکھو۔

\end{document}
